% ══════════════════════════════════════════════════════════════════
\chapter{Initialization Strategies}
\label{ch:initializers}
% ══════════════════════════════════════════════════════════════════

All initializers are defined for a weight matrix
$W \in \R^{d_{\mathrm{out}} \times d_{\mathrm{in}}}$
(the layout of \texttt{nn.Linear} in PyTorch, where each \emph{row}
is the weight vector of one output neuron).
We write $d = d_{\mathrm{in}}$ throughout (the number of input
connections per neuron).

% ──────────────────────────────────────────────────────────────────
\section{Baselines}

\noindent
These are standard initializers from the literature, providing
reference points for our analysis.

\subsection{He (Kaiming) initialization~\cite{he2015delving}}
\label{sec:init_he}

\begin{equation}\label{eq:he}
  W_{ij} \sim \mathcal{N}\!\left(0,\; \frac{2}{d}\right),
  \qquad\text{i.e.}\quad
  \sigma_W = \sqrt{\frac{2}{d}}.
\end{equation}
Derived by requiring $\E[(a^{(l)})^2] = \E[(a^{(l-1)})^2]$
under ReLU activation (see He et al., 2015, for the full derivation).

\subsection{Xavier (Glorot) initialization~\cite{glorot2010understanding}}

\begin{equation}
  W_{ij} \sim \mathcal{N}\!\left(0,\; \frac{2}{d_{\mathrm{in}} + d_{\mathrm{out}}}\right).
\end{equation}
Designed for linear/sigmoid/tanh activations; does not account for
ReLU's halving effect.  Included as a reference baseline.

\subsection{Uniform He}

\begin{equation}
  W_{ij} \sim \mathcal{U}(-a, a), \qquad a = \sqrt{\frac{6}{d}},
\end{equation}
so that $\Var(W_{ij}) = a^2/3 = 2/d$, matching He variance with a
uniform distribution.

\subsection{Orthogonal (generic)}

\begin{equation}
  W = g \cdot Q, \qquad Q \text{ from QR decomposition of }
  G \sim \mathcal{N}(0,1)^{d_{\mathrm{out}} \times d_{\mathrm{in}}}.
\end{equation}
The matrix $Q$ is obtained from the QR decomposition of a random
Gaussian matrix $G$, ensuring its rows form an orthonormal set.
The gain parameter $g$ scales all entries uniformly: $g = 1$ gives
unit-norm rows, $g = \sqrt{2}$ gives He-scaled rows
(see \texttt{orthogonal\_he} in Section~\ref{sec:orthogonal_family}).
This generic form is not used directly in our experiments ---
we use \texttt{orthogonal\_he} ($g = \sqrt{2}$) instead.
We define it here as the base construction.

% ──────────────────────────────────────────────────────────────────
\section{Row-Centered Family}

\noindent
Row centering was proposed to combat geometric collapse: by making
each neuron blind to the mean of its inputs, the hope was to reduce
the DC drift that accelerates angle contraction under the arc-cosine
kernel.  This family explores the consequences of that constraint.

\begin{definition}[Row centering]
\label{def:row_centering}
Given a weight matrix $W$, the \emph{row-centered} version is
\[
  \tilde{W}_{ij} = W_{ij} - \frac{1}{d}\sum_{k=1}^{d} W_{ik},
\]
so that each row sums to zero: $\sum_j \tilde{W}_{ij} = 0$ for all~$i$.
Here $d = d_{\mathrm{in}}$ is the number of input connections per neuron
(i.e., the length of each row).
\end{definition}

\textbf{Motivation.}  When row sums are zero, the neuron output
$z_i = \sum_j W_{ij} a_j$ can be rewritten as follows.
Split each activation into its mean plus deviation:
$a_j = \bar{a} + (a_j - \bar{a})$, where
$\bar{a} = \frac{1}{d}\sum_j a_j$ is the mean activation.
Then:
\begin{align}
  z_i &= \sum_j W_{ij}\, a_j \notag \\
      &= \sum_j W_{ij}\,\bar{a}
         + \sum_j W_{ij}\,(a_j - \bar{a}) \notag \\
      &= \bar{a}\,\underbrace{\sum_j W_{ij}}_{= \,0\;\text{(row centering)}}
         + \sum_j W_{ij}\,(a_j - \bar{a}) \notag \\
      &= \sum_j W_{ij}\,(a_j - \bar{a}).
  \label{eq:rc_motivation}
\end{align}
The key point: we are \emph{not} adding $-\bar{a}$ by hand.
We are splitting $a_j$ into mean plus deviation, and the mean term
vanishes automatically because the row sums to zero.
The result is that the neuron only sees the \emph{deviations} from
the mean, not the mean itself.

\subsection{\texttt{row\_centered\_he}}

\begin{equation}
  W_{ij}^{(\mathrm{He})} \sim \mathcal{N}(0, 2/d),
  \qquad
  W_{ij} = W_{ij}^{(\mathrm{He})} - \frac{1}{d}\sum_{k} W_{ik}^{(\mathrm{He})}.
\end{equation}
No variance correction.  Centering reduces the per-element variance by
a factor of $(1 - 1/d)$:

\begin{proof}[Proof of the variance reduction]
Let $w_j \sim \mathcal{N}(0, \sigma^2)$ i.i.d.\ for $j = 1,\dots,d$.
After centering: $\tilde{w}_j = w_j - \bar{w}$, where
$\bar{w} = \frac{1}{d}\sum_k w_k$.  Then:
\begin{align*}
  \Var(\tilde{w}_j)
  &= \Var(w_j - \bar{w})
  = \Var(w_j) + \Var(\bar{w}) - 2\,\Cov(w_j, \bar{w}).
\end{align*}
Now $\Var(\bar{w}) = \sigma^2/d$ (variance of the mean of $d$
i.i.d.\ variables).  And:
\[
  \Cov(w_j, \bar{w})
  = \Cov\!\left(w_j,\; \frac{1}{d}\sum_{k=1}^{d} w_k\right)
  = \frac{1}{d}\,\Var(w_j)
  = \frac{\sigma^2}{d},
\]
since $w_j$ is independent of all $w_k$ with $k \neq j$.
Therefore:
\[
  \Var(\tilde{w}_j)
  = \sigma^2 + \frac{\sigma^2}{d} - 2\,\frac{\sigma^2}{d}
  = \sigma^2\!\left(1 - \frac{1}{d}\right).
  \qedhere
\]
\end{proof}

\noindent
Geometrically, centering projects each row onto the hyperplane
orthogonal to $\bm{1}$, removing one degree of freedom out of $d$.

\subsection{\texttt{row\_centered\_he\_var\_adj}}
\label{sec:rc_var_adj}

\begin{enumerate}
  \item Sample $W^{(\mathrm{He})} \sim \mathcal{N}(0, 2/d)$.
  \item Center each row: $\tilde{\bm{w}}_i = \bm{w}_i - \bar{w}_i \bm{1}$.
  \item Rescale each row to restore He variance:
        $\bm{w}_i \leftarrow \tilde{\bm{w}}_i \cdot
        \sqrt{2/d}\,/\,\operatorname{std}(\tilde{\bm{w}}_i)$.
\end{enumerate}

\textbf{Why the rescaling (step~3).}
After centering, each row has per-element standard deviation
$\approx \sigma_W\sqrt{1 - 1/d}$ (from the variance reduction above).
For $d = 784$, this is $\sigma_W \cdot 0.9994$ --- nearly identical,
but the rescaling ensures \emph{exactly} $\Var(W_{ij}) = 2/d$
(matching He).  Each row is divided by its empirical standard deviation
and multiplied by $\sqrt{2/d}$.
The result: centered rows (sum to zero) with exactly He variance.

\subsection{\texttt{row\_centered\_forward\_balanced} (diagnostic)}
\label{sec:rc_fwd_balanced}

\textbf{Motivation.}
In Chapter~\ref{ch:gradient_trap}, we show that row centering causes
a forward gain of $\sqrt{(\pi-1)/\pi} \approx 0.826$ per layer.
A natural question: \emph{can we simply increase the weight variance
to compensate?}  This initializer answers that question.

\textbf{Derivation.}
For a row-centered layer with variance $\Var(W_{ij}) = v$, the forward
gain is (from Chapter~\ref{ch:gradient_trap}):
\[
  g_{\mathrm{fwd}} = \sqrt{v \cdot \frac{d}{2} \cdot \frac{\pi-1}{\pi}}.
\]
Setting $g_{\mathrm{fwd}} = 1$ and solving for $v$:
\[
  v \cdot \frac{d}{2} \cdot \frac{\pi-1}{\pi} = 1
  \quad\Longrightarrow\quad
  v = \frac{2\pi}{(\pi-1)\,d} \approx \frac{2.934}{d}.
\]
So the required weight variance is:
\begin{equation}\label{eq:fwd_balanced_var}
  \Var(W_{ij}) = \frac{2\pi}{(\pi-1)\,d}.
\end{equation}

\textbf{Result.}
This is a \textbf{diagnostic} initializer: it proves that forcing
the forward gain to~$1.0$ pushes the backward gain to
$\sqrt{\pi/(\pi-1)} \approx 1.21$, causing backward explosion
($1.21^{50} \approx 1.3 \times 10^4$).
You cannot fix one direction without breaking the other
(Section~\ref{sec:fwd_bwd_coupling}).

\subsection{\texttt{row\_centered\_layer\_balanced}}
\label{sec:layer_balanced}

\textbf{Motivation.}
Instead of using the same weight variance at every layer, we use a
\emph{per-layer variance schedule} $\sigma_l$ that partially
compensates for the cumulative forward decay.

\textbf{Definition.}
\begin{equation}\label{eq:layer_balanced}
  \sigma_l = \sigma_{\mathrm{base}} \cdot r^{\,\eta\,(l - (L+1)/2)},
  \qquad
  \sigma_{\mathrm{base}} = \sqrt{\frac{2}{d}} \cdot \sqrt{\frac{\pi}{\pi-1}},
\end{equation}
where $r = \sqrt{(\pi-1)/\pi}$, $l \in \{1,\dots,L\}$ is the layer
index, $L$ is the total number of layers, and $\eta \in [0,1]$ controls
the correction strength.

Each layer's weight matrix is sampled $W \sim \mathcal{N}(0, \sigma_l)$,
then row-centered and rescaled to have per-row standard deviation
$\sigma_l$.

\textbf{Intuition.}
The parameter $\eta$ controls how much per-layer compensation to apply:
\begin{itemize}
  \item $\eta = 0$: uniform variance across layers, same as
        \texttt{forward\_balanced} (forward gain~$\approx 1.0$,
        backward gain~$\approx 1.21$ at every layer).
  \item $\eta = 1$: full per-layer correction, with early layers
        getting smaller variance and later layers getting larger
        variance to fight cumulative decay.
  \item Intermediate $\eta$: a partial correction.
\end{itemize}

\begin{flagbox}
  The base standard deviation $\sigma_{\mathrm{base}}$ includes the
  $\sqrt{\pi/(\pi-1)}$ factor, which is the \textbf{forward-balanced}
  correction (Section~\ref{sec:rc_fwd_balanced}), not the variance-adjusted
  base.  Consequently, at $\eta = 0$ this initializer gives:
  \begin{itemize}
    \item Forward gain $\approx 1.0$ (not~$0.826$)
    \item Backward gain $\approx 1.21$ (not~$1.0$)
  \end{itemize}
  This differs from \texttt{row\_centered\_he\_var\_adj}, which at the
  same depth would give forward gain~$\approx 0.826$ and backward
  gain~$\approx 1.0$.  The distinction matters when interpreting the
  $\eta$-sweep results in Section~\ref{sec:eta_sweep}.
\end{flagbox}

The full mathematical derivation of why this scheme partially equalizes
gradients is given in Section~\ref{sec:eta_sweep}.

% ──────────────────────────────────────────────────────────────────
\section{Partial Centering}

\noindent
Since full row centering creates a gradient trap
(Chapter~\ref{ch:gradient_trap}), a natural follow-up:
\emph{what if we only partially center?}
This family interpolates between He (no centering, good gradients)
and full row centering (geometry-motivated, but trapped gradients).

\subsection{\texttt{partial\_centered\_he}}

\begin{equation}\label{eq:partial}
  W_{ij} = W_{ij}^{(\mathrm{He})}
           - \alpha \cdot \frac{1}{d}\sum_{k} W_{ik}^{(\mathrm{He})},
\end{equation}
followed by rescaling each row to restore He variance.

The parameter $\alpha \in [0,1]$ interpolates between pure He
($\alpha = 0$, no centering) and full row centering ($\alpha = 1$).
At intermediate $\alpha$, the row sums are reduced but not zero:
$\sum_j W_{ij} = (1-\alpha)\sum_j W_{ij}^{(\mathrm{He})}$.

Default: $\alpha = 0.5$.

% ──────────────────────────────────────────────────────────────────
\section{Orthogonal Family}
\label{sec:orthogonal_family}

\noindent
Orthogonal initialization preserves norms exactly in the linear part
of the forward pass (before the nonlinearity).  Combined with He
scaling, it provides the benefits of structured weights without the
zero-sum constraint that causes the gradient trap.

\subsection{\texttt{orthogonal\_he}}

\begin{equation}
  W = \sqrt{2} \cdot Q, \qquad Q \text{ orthogonal from QR}.
\end{equation}
The gain $\sqrt{2}$ matches He scaling: since ReLU zeroes ${\sim}50\%$
of outputs, the factor of~2 compensates for the halving.
Orthogonal rows are linearly independent but do \textbf{not} have
the zero-sum constraint, so there is no gradient trap.

\subsection{\texttt{orthogonal\_tuned}}

\begin{equation}
  W = \sqrt{1.65} \cdot Q.
\end{equation}
Uses the tuned factor $1.65$ instead of~$2$.

In our experiments (Section~\ref{sec:orthogonal_analysis}),
\texttt{orthogonal\_tuned} produces \textbf{identical results} to
\texttt{orthogonal\_he} on all scale-invariant metrics ($k$-NN
accuracy, distance correlation).  This is because for square matrices
with the same random seed, the QR decomposition produces the same
orthogonal matrix $Q$; only the scalar gain differs, which does not
affect angle-based or rank-based metrics.
We include it for completeness but focus on \texttt{orthogonal\_he}
in our analysis.

% ──────────────────────────────────────────────────────────────────
\section{Experimental Initializers}

\noindent
These are exploratory approaches that test specific hypotheses about
fixing the gradient trap or the kernel distortion.

\subsection{\texttt{centered\_with\_dc\_he}}

\begin{enumerate}
  \item Sample and center as in \texttt{row\_centered\_he\_var\_adj}.
  \item Add a constant DC offset to every element:
        \[
          W_{ij} \leftarrow W_{ij} + \delta \cdot \sqrt{2/d},
          \qquad \delta = 0.1 \text{ (default)}.
        \]
\end{enumerate}
This breaks the zero-sum constraint
($\sum_j W_{ij} = d \cdot \delta \sqrt{2/d} \neq 0$)
while preserving most of the centering benefit.

\subsection{\texttt{kernel\_preserving}}
\label{sec:kp_init}

This initializer directly minimizes the kernel distortion via
optimization.  Starting from He initialization, it runs 200~Adam steps
(learning rate 0.01) on the objective:
\begin{equation}\label{eq:kp_objective}
  \min_W \;
  \sum_{i,j}
  \left(
    \frac{\langle \relu(W\bm{x}_i),\, \relu(W\bm{x}_j)\rangle}{d_{\mathrm{out}}}
    - \langle \bm{x}_i, \bm{x}_j \rangle
  \right)^{\!2}
  + \lambda
  \sum_i
  \left(
    \frac{\|\relu(W\bm{x}_i)\|^2}{d_{\mathrm{out}}} - 1
  \right)^{\!2},
\end{equation}
where $\{\bm{x}_i\}$ are 500 random unit vectors (data-independent
proxy) and $\lambda = 1.0$.

\begin{flagbox}
  \textbf{Normalization artifact.}
  Both terms in~\eqref{eq:kp_objective} are divided by
  $d_{\mathrm{out}}$.  The norm constraint therefore asks for
  $\|\relu(W\bm{x})\|^2 / d_{\mathrm{out}} \approx 1$, i.e.,
  the optimizer targets output vectors of length
  $\|\relu(W\bm{x})\| \approx \sqrt{d_{\mathrm{out}}}$.
  For our layer width $d = 784$, this means
  $\sqrt{784} = 28$, not~$1$.

  When we later observe (Section~\ref{sec:kp_analysis}) that the
  optimized weights are ``$11.5\times$ inflated'' compared to He,
  part of this inflation is simply because the optimizer is targeting
  larger output norms than He does --- it is partly a normalization
  artifact, not necessarily a meaningful structural difference.
\end{flagbox}

\subsection{\texttt{custom\_variance}}

\begin{equation}
  W_{ij} \sim \mathcal{N}(\mu, \sigma^2),
\end{equation}
with user-specified $\mu$ (default~0) and $\sigma^2$ (default $2/d$).
Used for variance sweep experiments.

% ──────────────────────────────────────────────────────────────────
\section{Summary Table}

\begin{table}[h]
\centering
\small
\begin{tabular}{lcccc}
\toprule
Initializer & Row sum $=0$? & Var adjusted? & Gradient trap? & Best for \\
\midrule
\texttt{he}                     & No      & N/A       & No      & Gradient baseline \\
\texttt{xavier}                 & No      & N/A       & No      & Linear/tanh \\
\texttt{uniform\_he}            & No      & N/A       & No      & Uniform baseline \\
\midrule
\texttt{row\_centered\_he}      & Yes     & No        & Yes     & Geometry baseline \\
\texttt{rc\_he\_var\_adj}       & Yes     & Yes       & Yes     & Geom.\ + correct var \\
\texttt{rc\_fwd\_balanced}      & Yes     & Yes (2.93)& Yes     & Diagnostic only \\
\texttt{rc\_layer\_balanced}    & Yes     & Per-layer & Yes     & Geom.\ + grad.\ balance \\
\midrule
\texttt{partial\_centered\_he}  & Partial & Yes       & Partial & Trade-off \\
\midrule
\texttt{orthogonal\_he}         & No      & N/A       & No      & Geom.\ without trap \\
\texttt{orthogonal\_tuned}      & No      & N/A       & No      & $\equiv$ orthogonal\_he \\
\midrule
\texttt{centered\_with\_dc}     & Nearly  & Yes       & Partial & Breaking the trap \\
\texttt{kernel\_preserving}     & No      & Optimized & No      & Shallow geom. \\
\texttt{custom\_variance}       & No      & Custom    & No      & Experiments \\
\bottomrule
\end{tabular}
\caption{Summary of all initialization strategies.}
\label{tab:init_summary}
\end{table}

\begin{remark}[Key initializers for discussion]
The most important rows in this table are:
\begin{itemize}
  \item \textbf{\texttt{he}}: the standard baseline --- good gradients,
        reasonable geometry.
  \item \textbf{\texttt{rc\_he\_var\_adj}}: the advisor's proposal
        with variance correction --- this is what we analyze in depth.
  \item \textbf{\texttt{rc\_fwd\_balanced}}: diagnostic variant that
        reveals the gradient trap coupling (Section~\ref{sec:fwd_bwd_coupling}).
  \item \textbf{\texttt{orthogonal\_he}}: the practical winner ---
        best geometry with no gradient trap
        (Section~\ref{sec:orthogonal_analysis}).
  \item \textbf{\texttt{kernel\_preserving}}: shows that direct
        optimization finds the \emph{opposite} of centering
        (Section~\ref{sec:kp_analysis}).
\end{itemize}
The remaining initializers are supporting experiments that help us
understand the landscape of possible solutions.
\end{remark}
